\begin{englishabstract}
Water-to-air direct wireless optical communication allows underwater nodes to exchange data with aerial platforms without a surface relay. When an upward beam leaves seawater, the bare water--air interface totally reflects light beyond $48.75^\circ$ and refracts sub-critical oblique beams toward grazing directions; wind waves roughen the surface and move any floating platform, while position jitter, pointing error and turbulence fading further degrade link reliability. This thesis places a transmissive liquid-crystal reconfigurable intelligent surface (RIS) on the sea surface, mounted in contact with the water, which applies a programmable per-tile phase before the beam enters air. System models are built for the coordinate frames and device, platform motion and phase design, the incident and receiver-side complex fields, and link performance over a dynamic sea, and the performance gain provided by the RIS is evaluated.

First, a unified set of coordinate frames is established. For a homogeneous water body, a Gaussian source and an absorbing RIS boundary, the radiative transfer equation is solved by a collision series that converges with the single-scattering albedo as its ratio, and layered Fermat paths, area Jacobians and Fresnel transmission are combined into a complex-field transfer relation. A random sea surface is generated from the Neumann directional spectrum. A $1\,\mathrm m\times1\,\mathrm m$ glass-laminated liquid-crystal RIS with $40\times40$ tiles is designed; a 1\,$\mu$m electrode pitch gives a single-axis steering limit of $15.426^\circ$ and a command step of $0.021^\circ$, and water-contact mounting postpones total internal reflection to $72.15^\circ$.

Second, the buoyancy frame is trimmed and a six-degree-of-freedom rigid-body model coupling thin-plate radiation/diffraction loads, an analytic side-wall yaw moment and a passive radiation memory gives RIS attitudes and trajectories for wind speeds of 3--12\,m/s. For a given pose, the receiver is placed inside the direction cone reachable by all tiles, each tile is assigned a flat-top or robust target landing point from whitened radial energy quantiles, and loadable phases are obtained by layered Fermat-path inversion, clamping, quantization and piston calibration. A phase-control scheme that combines pose prediction with robust spots is proposed.

Third, the irradiance on the RIS lower face is decomposed into a ballistic part and a full Henyey--Greenstein scattered part; over 18 frozen snapshots the differences from independent Monte Carlo results match the Monte Carlo noise. For any loadable phase, the receiver-side complex field is computed with per-tile diffraction integrals, the chirp-$z$ transform, power-conserving deposition and pixel-integrated coherent superposition; over 228 phases at 10--50\,m and 3--12\,m/s its received power agrees with independent Monte Carlo within 10\%.

Finally, on a frozen dynamic rough sea, a no-RIS direct link computed by Monte Carlo is compared with four RIS schemes, Steering-only, Flat top, Robust ABER/OP and Max mean power, all sharing one jittered receiver. At receiver heights of 30\,m and above, the direct link cannot reach an average bit error rate of $10^{-3}$ within 90\,dBm, whereas the best RIS scheme needs about 42\,dBm, almost independent of height and wind speed. Robust ABER/OP and Flat top are the two clearly best methods, requiring about 6.5\,dB less power than Steering-only and about 14\,dB less than the maximum-mean-power design at 30\,m, and the robust spot further leads in outage probability under strong oceanic turbulence. The models and results support device selection, phase design and link evaluation for RIS-assisted water-to-air optical communication.

\englishkeyword{Water-to-Air Optical Wireless Communication, Reconfigurable Intelligent Surface, Liquid-Crystal Phase Modulation, Six-Degree-of-Freedom Motion, Complex Field, Robust Beam Shaping}
\end{englishabstract}
